\chapter{Língua Inglesa}
\label{cap:ingles-teoria}

Este capítulo trata a prova de Inglês pelo que ela realmente é: um teste de
\textbf{leitura instrumental} sobre texto técnico/jornalístico curto, não de
conversação nem de gramática decorada isoladamente. Cada seção cobre um
mecanismo da língua que a interpretação de texto explora --- referência
pronominal, conectivos, modais, vocabulário em contexto --- com o raciocínio
para resolver, não só a lista de palavras. A calibração de estilo (que tipo
de texto, que tipo de pegadinha) está na \texttt{apostila/}.

\section{Como ler para prova: \textit{skimming} e \textit{scanning}}

\begin{conceito}[Duas técnicas de leitura, para dois tipos de pergunta]
\begin{itemize}
\item \textbf{\textit{Skimming}}: leitura \textbf{rápida e superficial} do
texto inteiro, para captar a ideia geral, o gênero e o propósito --- sem
parar em palavra desconhecida. Use no primeiro contato com o texto, antes
de olhar as perguntas, e para responder questões de “ideia principal” ou
“propósito do texto”.
\item \textbf{\textit{Scanning}}: busca \textbf{dirigida} por uma
informação específica (um nome, uma data, uma palavra) --- os olhos varrem
o texto procurando o ponto exato, ignorando o resto. Use depois de ler a
pergunta, quando ela pede um detalhe pontual (referência de pronome,
sentido de uma palavra específica, se uma afirmativa bate com o texto).
\end{itemize}
\end{conceito}

\begin{regrapratica}[A ordem que economiza tempo de prova]
Não tente entender \textbf{cada palavra} do texto antes de ir às perguntas
--- é a armadilha mais comum de quem trata a prova como teste de
vocabulário. Faça um \textit{skimming} rápido (ideia geral + gênero), \emph{
depois} leia a primeira pergunta, e só então volte ao texto com
\textit{scanning} para achar o trecho exato que ela cobra. A resposta de
questão de detalhe está sempre em uma frase ou duas, localizáveis --- não
exige reler o texto inteiro de novo.
\end{regrapratica}

\section{Referência pronominal}

\begin{conceito}[O pronome aponta para trás --- em regra, para o mais próximo relevante]
Pronomes (\textit{it, this, that, they, her, his, these, those}) existem
para \textbf{evitar repetir} um substantivo já mencionado --- exatamente
como a anáfora do português (Capítulo \ref{cap:portugues-teoria}). Resolver
a referência é achar, \textbf{voltando ao texto}, a que substantivo/ideia o
pronome se refere. A armadilha: a prova costuma colocar \textbf{dois}
substantivos candidatos na mesma vizinhança do pronome, e um deles é
gramaticalmente plausível mas semanticamente errado.
\end{conceito}

\begin{exemplo}[Resolvendo um pronome ambíguo à primeira vista]
\textit{“The professor asked the student to review the report before her
presentation.”}

“Her” pode, à primeira vista, parecer se referir a “the student” (mais
próximo na frase) --- mas o sentido completo depende de \textbf{quem vai
apresentar}: se o contexto no texto deixa claro que é a aluna quem vai
apresentar, “her” = the student; se é a professora quem vai apresentar
depois de receber a revisão, “her” = the professor. \textbf{Proximidade
sintática não decide sozinha} --- é preciso o sentido da frase completa,
às vezes até do parágrafo anterior, para confirmar.
\end{exemplo}

\begin{cuidado}[O distrator do pronome usa o candidato mais óbvio, não o correto]
Quando a prova pergunta a que um pronome se refere, uma das alternativas
erradas quase sempre é o substantivo \textbf{fisicamente mais próximo} do
pronome no texto --- porque parece a resposta “óbvia” para quem não voltou
a conferir. A técnica correta é sempre \textbf{retornar à frase} (às vezes
à anterior) e testar a substituição: troque o pronome pelo candidato e
veja se a frase mantém o sentido lógico do texto inteiro, não só da oração
isolada.
\end{cuidado}

\section{Conectivos e o valor lógico que carregam}

\begin{conceito}[Cada conectivo sinaliza uma relação lógica específica]
Assim como em português (Capítulo \ref{cap:portugues-teoria},
Seção~\ref{sec:correlacao}), a interpretação em inglês cobra menos “que
classe gramatical é essa palavra” e mais “\textbf{que relação lógica} ela
introduz entre as duas partes do texto”:

\begin{center}
\begin{tabular}{@{}p{3.2cm}p{4.4cm}p{4cm}@{}}
\toprule
\textbf{Relação} & \textbf{Conectivos} & \textbf{Sinal} \\
\midrule
Consequência/efeito & thus, therefore, hence, so, as a result & o que vem
depois é \textbf{resultado} do que veio antes \\
Contraste/concessão & however, whereas, but, although, despite, yet & o
que vem depois \textbf{contradiz ou relativiza} o que veio antes \\
Adição & moreover, furthermore, in addition, besides, also & reforça no
\textbf{mesmo sentido} do que já foi dito \\
Causa & because, since, as, due to & o que vem depois é o \textbf{motivo}
do que veio antes \\
Condição & if, unless, provided that & \textit{unless} é
\textbf{condição negativa} (= “if not”) \\
Tempo & meanwhile, while, when, after, before & situa \textbf{quando} algo
acontece em relação a outra ação \\
Exemplificação & for example, for instance, such as & introduz um
\textbf{caso concreto} do que foi dito de forma geral \\
\bottomrule
\end{tabular}
\end{center}
\end{conceito}

\begin{exemplo}[Trocando o conectivo certo pela relação certa]
\textit{“The company's servers were outdated. \textbf{Thus}, the migration
to cloud infrastructure became urgent.”}

“Thus” marca que a migração urgente é \textbf{consequência} de os
servidores estarem desatualizados --- não é um exemplo, nem uma condição,
nem um contraste. Se o texto trocasse “thus” por “however”, o sentido
mudaria para uma \textbf{oposição} --- e a frase deixaria de fazer sentido
lógico com o contexto dado.
\end{exemplo}

\begin{cuidado}[A banca confunde causa/consequência com contraste]
O erro mais comum em questão de conectivo é a alternativa trocar a relação
\textbf{real} do texto (efeito, adição, causa) por \textbf{contraste}, ou
vice-versa --- porque são categorias que soam parecidas quando lidas
rápido demais. Antes de responder, releia a frase completa (a que vem
\emph{antes} e a que vem \emph{depois} do conectivo) e pergunte: a segunda
ideia \textbf{decorre} da primeira, \textbf{se opõe} a ela, ou só a
\textbf{reforça}?
\end{cuidado}

\section{Verbos modais}

\begin{conceito}[O que cada modal sinaliza sobre a atitude do autor]
Modal não é só gramática de conjugação --- cada um carrega um
\textbf{grau de certeza ou obrigação} que muda a leitura do texto:

\begin{center}
\begin{tabular}{@{}p{2.4cm}p{4.6cm}p{4.4cm}@{}}
\toprule
\textbf{Modal} & \textbf{Sentido} & \textbf{Força} \\
\midrule
must & obrigação/necessidade forte, ou dedução quase certa & alta \\
should/ought to & sugestão, conselho, expectativa & média \\
have to & obrigação (geralmente externa, imposta por regra/circunstância) &
alta \\
can/could & possibilidade, permissão, capacidade (\textit{could} é mais
formal/hipotético) & variável \\
may/might & possibilidade --- \textit{might} indica \textbf{menor}
certeza que \textit{may} & baixa a média \\
\bottomrule
\end{tabular}
\end{center}

O foco de prova não é decorar a tabela de regras de uso do modal em
isolado, é reconhecer \textbf{que atitude o autor do texto está expressando}
quando o usa --- um “should” no meio de um parágrafo é uma
\textbf{recomendação}, não um fato relatado como certo.
\end{conceito}

\begin{exemplo}[O modal muda a força da afirmação]
\textit{“Developers should review AI-generated code before merging it.”}
--- o autor está \textbf{recomendando} uma prática, não afirmando que é
uma regra universal seguida por todos.

\textit{“Developers must review AI-generated code before merging it,
according to the company's new policy.”} --- aqui é \textbf{obrigação}
formal, imposta por uma política --- uma leitura mais forte que a
anterior, mesmo a frase sendo quase idêntica.
\end{exemplo}

\section{Vocabulário em contexto}

\begin{conceito}[A mesma palavra muda de classe --- e de sentido --- conforme a função]
Assim como em português (Capítulo \ref{cap:portugues-teoria}, Seção de
Classes de Palavras), várias palavras do inglês técnico funcionam ora como
\textbf{substantivo}, ora como \textbf{verbo}, e o sentido muda conforme a
posição na frase --- não existe glossário fixo que resolva sozinho, é
preciso olhar a função sintática:

\begin{itemize}
\item \textbf{browse}: verbo (“to browse the web” = navegar) ou, mais raro,
substantivo informal.
\item \textbf{rate}: verbo (“to rate a product” = avaliar) ou substantivo
(“a high failure rate” = uma taxa).
\item \textbf{access}: verbo (“to access a system” = acessar) ou
substantivo (“access to the system” = o acesso).
\item \textbf{post}: verbo (“to post an update” = publicar) ou substantivo
(“read the post” = a publicação).
\end{itemize}
\end{conceito}

\begin{regrapratica}[Como decidir a classe de uma palavra ambígua]
Olhe a \textbf{posição} na frase: logo depois de \textit{to}, de um modal
(\textit{can, should}), ou concordando como núcleo de uma oração $\to$
provavelmente \textbf{verbo}. Precedida de artigo (\textit{a, the}) ou de
adjetivo, ou seguida de preposição que a modifica $\to$ provavelmente
\textbf{substantivo}. O mesmo teste de substituição usado em português
funciona aqui: tente trocar por um verbo/substantivo inequívoco e veja se a
frase continua de pé com a mesma estrutura.
\end{regrapratica}

\begin{cuidado}[Falsos cognatos --- parece português, não é]
Palavras parecidas com o português enganam quando o sentido diverge:
\textit{“actually”} não é “atualmente” (é “na verdade”);
\textit{“eventually”} não é “eventualmente” (é “finalmente,
eventualmente no sentido de resultado final”); \textit{“pretend”} não é
“pretender” (é “fingir”); \textit{“parents”} não é “parentes” (é
“pais”). Em texto técnico de TI, cuidado especial com \textit{“library”}
(biblioteca de código, não “livraria”) e \textit{“realize”} (perceber, dar-se
conta --- não necessariamente “realizar/executar”).
\end{cuidado}

\section{Compreensão global: ideia principal, propósito e gênero}

\begin{conceito}[Três perguntas diferentes sobre o texto como um todo]
\begin{itemize}
\item \textbf{Ideia principal} (\textit{main idea}): o resumo do que o
texto \textbf{diz}, em uma frase --- não um detalhe isolado de um
parágrafo, mas o fio condutor de todos eles.
\item \textbf{Propósito} (\textit{purpose}): \textbf{por que} o autor
escreveu --- para informar, persuadir, alertar, vender, instruir. Um
mesmo assunto pode ser tratado com propósitos diferentes (um texto sobre
IA pode informar sobre uma pesquisa ou alertar sobre um risco).
\item \textbf{Gênero} (\textit{genre}): \textbf{que tipo} de texto é ---
anúncio, resenha, \textit{abstract} acadêmico, notícia, post de blog de
segurança. Reconhecer o gênero ajuda a prever o propósito (um
\textit{abstract} acadêmico tem propósito de relatar achado de pesquisa; um
anúncio tem propósito de vender).
\end{itemize}
\end{conceito}

\begin{regrapratica}[Onde procurar a ideia principal]
Em textos curtos de prova, a ideia principal costuma estar \textbf{no
primeiro parágrafo} (a tese/tópico) e ser \textbf{retomada} no último. Se a
pergunta pede “o principal ponto do texto”, desconfie de alternativas que
citam um \textbf{detalhe} de um único parágrafo do meio --- detalhe real
não é a mesma coisa que ideia central, mesmo estando literalmente no
texto.
\end{regrapratica}

\section{Questões de julgamento (afirmativas verdadeiro/falso)}

\begin{conceito}[Cada afirmativa é um teste independente contra o texto]
O formato “julgue as afirmativas I, II, III e marque a alternativa que
indica as corretas” exige comparar \textbf{cada} afirmativa,
\textbf{isoladamente}, contra o que o texto \textbf{efetivamente} diz --- o
mesmo princípio da distinção entre informação explícita e extrapolação do
capítulo de Português (Capítulo \ref{cap:portugues-teoria}). Uma
afirmativa é falsa se: (1) contradiz o texto diretamente; (2) inverte um
detalhe (troca um número, um nome, uma condição); ou (3) acrescenta uma
informação que o texto \textbf{não} sustenta, mesmo que pareça plausível no
mundo real.
\end{conceito}

\begin{exemplo}[Os três jeitos de uma afirmativa estar errada]
Texto-base: \textit{“The report, based on a survey of 500 developers,
found that AI coding assistants reduced debugging time by 20\% on
average.”}

\begin{itemize}
\item Contradição direta: \textit{“The report found that AI assistants
increased debugging time.”} --- inverte o resultado.
\item Detalhe trocado: \textit{“The survey included 5,000 developers.”}
--- trocou o número (500 por 5.000).
\item Extrapolação: \textit{“The report proves that all developers should
adopt AI assistants immediately.”} --- o texto relata um achado estatístico
médio, não uma recomendação universal e imediata; “proves” e “should...
immediately” vão além do que o texto sustenta.
\end{itemize}
\end{exemplo}

\begin{regrapratica}[Resolva uma afirmativa de cada vez]
Não tente julgar as três afirmativas simultaneamente nem “pelo padrão da
resposta”. Releia o texto \textbf{uma vez para cada afirmativa}, marque
verdadeiro/falso individualmente, e só então combine com as alternativas
de resposta (que geralmente cobrem combinações como “I e III” ou “apenas
II”). Julgar em bloco é o que leva a arrastar um erro de uma afirmativa
para a interpretação das outras.
\end{regrapratica}

\section{Roteiro de revisão do capítulo}

\begin{regrapratica}[Sequência para qualquer questão de interpretação em inglês]
\begin{enumerate}
\item \textit{Skimming} no texto inteiro: gênero, ideia geral, propósito.
\item Leia a pergunta antes de caçar a resposta no texto.
\item Pergunta de \textbf{pronome}? Volte à frase (ou à anterior) e teste
a substituição --- não confie no candidato mais próximo.
\item Pergunta de \textbf{conectivo}? Releia a oração antes e depois dele;
identifique a relação lógica real (causa, contraste, adição, condição).
\item Pergunta de \textbf{vocabulário}? Olhe a função sintática da palavra
na frase (verbo ou substantivo) antes de decidir o sentido.
\item Pergunta de \textbf{compreensão global}? Desconfie de alternativas
que citam só um detalhe isolado, ou que acrescentam informação não
sustentada pelo texto (extrapolação).
\item Afirmativas I/II/III? Julgue uma de cada vez, contra o texto, antes
de olhar as combinações de resposta.
\end{enumerate}
\end{regrapratica}
